% Options for packages loaded elsewhere
\PassOptionsToPackage{unicode}{hyperref}
\PassOptionsToPackage{hyphens}{url}
\PassOptionsToPackage{dvipsnames,svgnames,x11names}{xcolor}
%
\documentclass[
  letterpaper,
  DIV=11,
  numbers=noendperiod]{scrartcl}

\usepackage{amsmath,amssymb}
\usepackage{iftex}
\ifPDFTeX
  \usepackage[T1]{fontenc}
  \usepackage[utf8]{inputenc}
  \usepackage{textcomp} % provide euro and other symbols
\else % if luatex or xetex
  \usepackage{unicode-math}
  \defaultfontfeatures{Scale=MatchLowercase}
  \defaultfontfeatures[\rmfamily]{Ligatures=TeX,Scale=1}
\fi
\usepackage{lmodern}
\ifPDFTeX\else  
    % xetex/luatex font selection
\fi
% Use upquote if available, for straight quotes in verbatim environments
\IfFileExists{upquote.sty}{\usepackage{upquote}}{}
\IfFileExists{microtype.sty}{% use microtype if available
  \usepackage[]{microtype}
  \UseMicrotypeSet[protrusion]{basicmath} % disable protrusion for tt fonts
}{}
\makeatletter
\@ifundefined{KOMAClassName}{% if non-KOMA class
  \IfFileExists{parskip.sty}{%
    \usepackage{parskip}
  }{% else
    \setlength{\parindent}{0pt}
    \setlength{\parskip}{6pt plus 2pt minus 1pt}}
}{% if KOMA class
  \KOMAoptions{parskip=half}}
\makeatother
\usepackage{xcolor}
\setlength{\emergencystretch}{3em} % prevent overfull lines
\setcounter{secnumdepth}{-\maxdimen} % remove section numbering
% Make \paragraph and \subparagraph free-standing
\ifx\paragraph\undefined\else
  \let\oldparagraph\paragraph
  \renewcommand{\paragraph}[1]{\oldparagraph{#1}\mbox{}}
\fi
\ifx\subparagraph\undefined\else
  \let\oldsubparagraph\subparagraph
  \renewcommand{\subparagraph}[1]{\oldsubparagraph{#1}\mbox{}}
\fi


\providecommand{\tightlist}{%
  \setlength{\itemsep}{0pt}\setlength{\parskip}{0pt}}\usepackage{longtable,booktabs,array}
\usepackage{calc} % for calculating minipage widths
% Correct order of tables after \paragraph or \subparagraph
\usepackage{etoolbox}
\makeatletter
\patchcmd\longtable{\par}{\if@noskipsec\mbox{}\fi\par}{}{}
\makeatother
% Allow footnotes in longtable head/foot
\IfFileExists{footnotehyper.sty}{\usepackage{footnotehyper}}{\usepackage{footnote}}
\makesavenoteenv{longtable}
\usepackage{graphicx}
\makeatletter
\def\maxwidth{\ifdim\Gin@nat@width>\linewidth\linewidth\else\Gin@nat@width\fi}
\def\maxheight{\ifdim\Gin@nat@height>\textheight\textheight\else\Gin@nat@height\fi}
\makeatother
% Scale images if necessary, so that they will not overflow the page
% margins by default, and it is still possible to overwrite the defaults
% using explicit options in \includegraphics[width, height, ...]{}
\setkeys{Gin}{width=\maxwidth,height=\maxheight,keepaspectratio}
% Set default figure placement to htbp
\makeatletter
\def\fps@figure{htbp}
\makeatother

\KOMAoption{captions}{tableheading}
\makeatletter
\@ifpackageloaded{caption}{}{\usepackage{caption}}
\AtBeginDocument{%
\ifdefined\contentsname
  \renewcommand*\contentsname{Table of contents}
\else
  \newcommand\contentsname{Table of contents}
\fi
\ifdefined\listfigurename
  \renewcommand*\listfigurename{List of Figures}
\else
  \newcommand\listfigurename{List of Figures}
\fi
\ifdefined\listtablename
  \renewcommand*\listtablename{List of Tables}
\else
  \newcommand\listtablename{List of Tables}
\fi
\ifdefined\figurename
  \renewcommand*\figurename{Figure}
\else
  \newcommand\figurename{Figure}
\fi
\ifdefined\tablename
  \renewcommand*\tablename{Table}
\else
  \newcommand\tablename{Table}
\fi
}
\@ifpackageloaded{float}{}{\usepackage{float}}
\floatstyle{ruled}
\@ifundefined{c@chapter}{\newfloat{codelisting}{h}{lop}}{\newfloat{codelisting}{h}{lop}[chapter]}
\floatname{codelisting}{Listing}
\newcommand*\listoflistings{\listof{codelisting}{List of Listings}}
\makeatother
\makeatletter
\makeatother
\makeatletter
\@ifpackageloaded{caption}{}{\usepackage{caption}}
\@ifpackageloaded{subcaption}{}{\usepackage{subcaption}}
\makeatother
\makeatletter
\@ifpackageloaded{fontawesome5}{}{\usepackage{fontawesome5}}
\makeatother
\ifLuaTeX
  \usepackage{selnolig}  % disable illegal ligatures
\fi
\usepackage{bookmark}

\IfFileExists{xurl.sty}{\usepackage{xurl}}{} % add URL line breaks if available
\urlstyle{same} % disable monospaced font for URLs
\hypersetup{
  pdftitle={Course: Statistics for Social Justice},
  pdfauthor={Tyler George  stats\_tgeorge  stats-tgeorge},
  colorlinks=true,
  linkcolor={blue},
  filecolor={Maroon},
  citecolor={Blue},
  urlcolor={Blue},
  pdfcreator={LaTeX via pandoc}}

\title{Course: Statistics for Social Justice}
\usepackage{etoolbox}
\makeatletter
\providecommand{\subtitle}[1]{% add subtitle to \maketitle
  \apptocmd{\@title}{\par {\large #1 \par}}{}{}
}
\makeatother
\subtitle{eCOTS 2024}
\author{Tyler George \faIcon{x-twitter}
\href{https://twitter.com/stats_tgeorge}{stats\_tgeorge} \faIcon{github}
\href{https://github.com/stats-tgeorge}{stats-tgeorge}}
\date{}

\begin{document}
\maketitle

\subsection{Class Setting}\label{class-setting}

\begin{itemize}
\tightlist
\item
  One Course at a Time

  \begin{itemize}
  \tightlist
  \item
    18 class days including finals
  \item
    4 hours per day
  \item
    Students take, and faculty teach \ldots{} one course at a time
  \end{itemize}
\item
  11 students
\item
  all sophomores
\item
  no pre-requisites
\item
  Community Partners

  \begin{itemize}
  \tightlist
  \item
    Primary: \href{https://www.waypointservices.org/}{Waypoint}
  \item
    Secondary: \href{https://www.iowalegalaid.org/}{Iowa Legal Aid}
  \end{itemize}
\end{itemize}

\subsection{Collaborative classroom}\label{collaborative-classroom}

\includegraphics{../images/classroom.jpg}

\subsection{Course Creation Timeline}\label{course-creation-timeline}

\begin{itemize}
\item
  October 2022: Contacted Civic Engagement Office
\item
  November 2022: Met with Waypoint to Discuss Ideas
\item
  February 2023: Decided on Research Questions
\item
  July 2023: Iowa Legal Aid Joined
\item
  Early August 2023: Met to determine final asks and to discuss a
  community partner agreement
\item
  Mid-August 2023: Set course details
\item
  End August 2023: Received data, login information provided, and class
  started
\end{itemize}

\subsection{Course Organization}\label{course-organization}

\begin{itemize}
\tightlist
\item
  General outline

  \begin{itemize}
  \tightlist
  \item
    Engage with materials to understand disparity and inequity
  \item
    Training by community patterns to use their websites and data
  \item
    Descriptive statistics
  \item
    Start group project
  \end{itemize}
\item
  Throughout

  \begin{itemize}
  \tightlist
  \item
    Relevant talks
  \item
    Podcasts and videos
  \item
    Reflection and course discussions
  \end{itemize}
\end{itemize}

\subsection{Pedagogy}\label{pedagogy}

\begin{itemize}
\item
  Engage -\textgreater{} Reflect -\textgreater{} Group discussion
  followed by whole class discussion
\item
  Collaborative learning

  \begin{itemize}
  \tightlist
  \item
    The class was divided into 3 groups
  \item
    Each group then divided tasks
  \item
    TV on the wall to work on the same website
  \item
    Each day those tasks rotated
  \item
    Example: 1 - Controlling source website, 2 - typing record from the
    website, 3 - facilitating watching person typing to verify the
    record, 4 - working on report and/or presentation
  \end{itemize}
\item
  Presentations - students had to present progress on multiple occasions
  and the final project was presented to the community partner
\end{itemize}

\subsection{Unprecedented
Difficulties}\label{unprecedented-difficulties}

\begin{itemize}
\tightlist
\item
  Waypoint was great but I didn't receive the data and sign-in
  information until the week before class started.
\item
  Transportation
\item
  Only 2 sets of login information
\item
  Students had to sign a type of confidentiality agreement (not a
  formal)
\item
  Picking software for data visualization
\item
  Recording consistency
\end{itemize}

\subsection{Give Back}\label{give-back}

\begin{itemize}
\tightlist
\item
  Students gave a final presentation to partners
\item
  I found and fixed some student errors in data, then wrote metadata for
  all data files and turned it all back to Waypoint
\end{itemize}

\section{The End - Thank you!}\label{the-end---thank-you}



\end{document}
